%% NMM UED + Machine Learning paper
%% Style: article (JCP-compatible; switch to revtex4-2 before final submission)
%% Compile: pdflatex -> bibtex -> pdflatex -> pdflatex

\documentclass[12pt]{article}

\usepackage{newtxtext}
\usepackage[noamssymbols]{newtxmath}
\usepackage{graphicx}
\usepackage{caption}
\usepackage{soul}
\usepackage[letterpaper,margin=1in]{geometry}
\usepackage{amsmath}
\usepackage{booktabs}
% siunitx removed
\usepackage{xcolor}
\usepackage{url}
\usepackage[numbers,super,sort&compress]{natbib}
\usepackage{hyperref}
\hypersetup{colorlinks=true,linkcolor=blue,citecolor=blue,urlcolor=blue}

\linespread{1.5}
\frenchspacing
\date{}
\renewcommand\refname{References}

\renewenvironment{abstract}{\quotation}{\endquotation}

% Short-hands
\newcommand{\dsm}{\Delta sM}
\newcommand{\AAng}{\,\text{\AA}}
\newcommand{\anginv}{\,\text{\AA}^{-1}}
\newcommand{\TBD}[1]{\textcolor{red}{\textbf{[#1]}}}

%%=============================================================
\begin{document}

\title{\bfseries\boldmath
Machine Learning Inversion of MeV Ultrafast Electron Diffraction Data
Reveals Structural Dynamics of Photoexcited N-Methylmorpholine}

\author{
  \TBD{Author list} \\[4pt]
  \small Key Laboratory for Laser Plasmas (Ministry of Education)
    and School of Physics and Astronomy, \\
  \small Tsung-Dao Lee Institute, Shanghai Jiao Tong University,
    Shanghai, China \\[4pt]
  \small $^\ast$Corresponding author: \texttt{dxiang@sjtu.edu.cn}
}

\maketitle

%%------------------------------------------------------------
\begin{abstract}
We present a machine learning framework for direct inversion of
time-resolved ultrafast electron diffraction (UED) signals and apply it
to reveal the structural dynamics of photoexcited N-methylmorpholine
(NMM) with MeV electrons. NMM is a six-membered nitrogen-oxygen
heterocycle whose Rydberg-state photochemistry---dominated by coherent
nitrogen umbrella-mode oscillations near 650~fs period---has been
established by gas-phase photoionization spectroscopy.\cite{Zhang2017,Beedle2017}
Here we provide the first direct structural measurements. A multi-scale
one-dimensional convolutional neural network (NMMRegressorV2,
2.1$\times$10$^6$ parameters) is trained on $8\times10^4$ synthetic difference
scattering signals $\dsm(s)$ computed with relativistic DPWA electron
scattering factors, with Gaussian noise matched to the experimentally
measured pre-time-zero noise level ($\sigma_{\rm noise}\approx 5.5$~normalized
units). The network directly maps each $\dsm(s)$ to seven
interatomic distances---three key distances (O$\cdots$N, O$\cdots$C5,
N--C5) and four ancillary distances to fixed anchors (N--C2, N--C4,
C5--C2, C5--C4)---enabling unique 3D reconstruction of the mobile
atoms via trilateration, without iterative fitting. Applied to 200
bootstrap-resampled experimental datasets at 45 pump-probe time delays
($-370$ to $+1606$~fs), it yields time-resolved distance trajectories
with 1$\sigma$ bootstrap uncertainty bands. The network achieves a mean
absolute error (MAE) of 0.097~\AA\ on the three key distances and
0.148~\AA\ overall on the 8,000-sample validation set. The O$\cdots$N distance increases by $+0.45$~\AA\ ($16\%$)
within $\sim$130~fs of photoexcitation, directly reflecting nitrogen
planarization (sp$^3 \to$ sp$^2$) in the Rydberg state. Trilateration
from the seven predicted distances yields full 3D reconstructions of
the N3 and C5 atom positions at all 45 time points (100\% success
rate).
This work demonstrates noise-matched deep-learning inversion as a
practical complement to traditional iterative structure refinement for
UED data of medium-sized organic molecules.
\end{abstract}

%%=============================================================
\section*{I.\; Introduction}
%%=============================================================

\noindent
Ultrafast electron diffraction (UED) and ultrafast X-ray scattering (UXS)
directly probe gas-phase molecular structural dynamics on femtosecond to
picosecond timescales through changes in the molecular interference
pattern of scattered radiation.\cite{Zewail2000,Yang2018,Ihee2001,Minitti2015}
MeV-UED instruments, based on radio-frequency photocathode guns, provide
temporal resolution below 100~fs and access the full reciprocal-space
range up to $s \sim 15\anginv$.\cite{Qi2020}

The primary observable in a UED pump-probe experiment is the
time-resolved difference modified scattering signal
\begin{equation}
  \dsm(s,t) = sM(s,t) - \overline{sM}(s,t<0),
  \label{eq:dsm}
\end{equation}
where $sM(s) = s\,I_M(s)/I_A(s)$ is the modified scattering intensity
and $s = 4\pi\sin(\theta/2)/\lambda_e$ is the momentum transfer.
Inverting \eqref{eq:dsm} to obtain interatomic distances is an ill-posed
problem traditionally approached by fitting a parameterized structural
model.\cite{Centurion2022} Machine learning (ML) regression, trained on
synthetic data, offers a non-iterative alternative that requires no prior
knowledge of the reaction coordinate and delivers inference in under a
millisecond per frame.\cite{Kirrander2021}

N-methylmorpholine (NMM, C$_5$H$_{11}$NO) is a saturated six-membered
ring with nitrogen and oxygen heteroatoms. Its nitrogen center is
pyramidal (sp$^3$) in the ground state; near-UV photoexcitation
($\sim$200--230~nm) accesses 3s and 3p Rydberg states and launches a
coherent wavepacket in the nitrogen umbrella inversion mode.\cite{Zhang2017}
This umbrella motion oscillates at $\sim 650$~fs period, with an initial
Franck-Condon displacement of about 17° from the pyramidal
geometry.\cite{Zhang2017} The coherent oscillations survive internal
conversion from 3p to 3s,\cite{Beedle2017} dephasing only on a
$\sim 750$~fs timescale as the wavepacket couples to the vibrational
bath.\cite{Zhang2017} This non-ergodic behavior makes NMM an important
model for photochemistry of nitrogen heterocycles relevant to
pharmaceuticals and biology.\cite{Nix2006}
Despite this spectroscopic characterization, \emph{direct structural
measurements of the umbrella amplitude in physical distance units} have
not been reported. Here, we extract time-resolved O$\cdots$N,
O$\cdots$C5, and N--C5 distances from experimental MeV-UED data via ML
inversion, providing the first atom-resolved view of the umbrella
dynamics of NMM.

%%=============================================================
\section*{II.\; Experimental Details}
%%=============================================================

\subsection*{A.\; MeV-UED Instrument}

Experiments were performed using the MeV-UED facility at Shanghai
Jiao Tong University.\cite{Qi2020} An ultrashort electron bunch at
3~MeV kinetic energy was generated by a photocathode radio-frequency
(RF) gun and longitudinally compressed by a double-bend achromat (DBA)
system to an estimated pulse duration of tens of femtoseconds, with
timing jitter below $\sim 20$~fs. The electron beam size at the gas
cell interaction region was \TBD{$\sim$150}~$\mu$m (FWHM).

NMM (purity $\geq 99\%$) was delivered as an effusive molecular beam
\TBD{[nozzle description]} at room temperature.

\subsection*{B.\; Pump Laser}

Photoexcitation was accomplished using $\sim$200~nm UV laser pulses
(wavelength: \TBD{X}~nm; pulse duration: \TBD{$\sim$100}~fs FWHM)
derived from \TBD{laser system}. The fluence at the interaction region
was \TBD{X}~mJ/cm$^2$. This wavelength accesses the lowest 3s Rydberg
state of NMM, consistent with prior spectroscopic studies.\cite{Zhang2017}

\subsection*{C.\; Data Acquisition}

Diffraction patterns were recorded at 45 pump-probe time delays
ranging from $-370.1$ to $+1605.6$~fs (see Table~S1 in the
Supplementary Information for the complete delay list). The nine
pre-pump frames (t00--t08, $t < 0$) serve as ground-state reference.
A total of 240 experimental runs were collected. Scattering profiles
$I(s)$ contain 637 data points with pixel size
$\Delta s = 0.0245\anginv$, covering $s = 0.012$--$15.6\anginv$.
Bootstrap resampling (200 resamples) was applied across runs to
estimate statistical uncertainties.

\subsection*{D.\; Preprocessing Pipeline}

Experimental $\dsm(s,t)$ signals were preprocessed as follows before
network inference:
\begin{enumerate}
  \item \textit{Reference subtraction and NaN repair}: Linear
    interpolation of missing values; subtraction of the mean pre-pump
    signal $\overline{sM}_{\rm ref}(s)$ (averaged over 200 bootstrap
    samples and all $t{<}0$ frames).
  \item \textit{Baseline correction}: For $s > 5.0\anginv$, a
    second-order polynomial is fitted and subtracted to remove slow
    instrumental drift.
  \item \textit{Range masking}: Signal zeroed outside
    $s \in [1.5,\,9.0]\anginv$ (below: inelastic background; above:
    SNR~$<$~1).
  \item \textit{Grid resampling}: Interpolated from 637 experimental
    points to the 681-point uniform model grid
    ($s = 0$--$15\anginv$).
  \item \textit{Amplitude calibration}: Scaled by $\alpha \approx 0.02$,
    determined by grid search over $\alpha \in [0.02,\,0.04]$ to
    minimize pre-pump distance deviation from the DFT equilibrium
    geometry.
\end{enumerate}

%%=============================================================
\section*{III.\; Machine Learning Framework}
%%=============================================================

\subsection*{A.\; Molecular Model and Structural Sampling}

The NMM ground-state heavy-atom geometry is taken from a
B3LYP/6-311++G** DFT optimization (Fig.~\ref{fig:structure}).
Seven heavy atoms define the skeleton: ring carbons C1, C2, C4, C6;
ring oxygen O7; ring nitrogen N3; and N-methyl carbon C5. Hydrogen
atoms are rigidly attached at their ground-state offsets from parent
heavy atoms (translational-rigidity approximation), contributing to
scattering but not forming independent degrees of freedom.

Five atoms (C1, C2, C4, C6, O7) are held fixed at ground-state
positions, consistent with the expectation that photoexcitation
primarily drives N3 and C5.\cite{Zhang2017} Two atoms are sampled
randomly:
\begin{itemize}
  \item \textbf{N3}: uniformly within a cylinder of radius
    $R_N = 2.0\AAng$ centered on the ground-state position, with
    off-plane displacement of $2.0\AAng$.
  \item \textbf{C5}: uniformly within a sphere of radius
    $R_{C5} = 3.0\AAng$ centered on the ground-state C5 position.
\end{itemize}
Samples with $r_{N\text{-}O} > r_{C5\text{-}O}$ are rejected; the
effective acceptance rate is $\approx$68\%.

\subsection*{B.\; Scattering Signal Computation}

Modified scattering signals are computed within the independent-atom
model (IAM):\cite{Centurion2022}
\begin{equation}
  sM(s) = s\sum_{i\neq j} f_i(s)\,f_j(s)\,
  \frac{\sin(s\,r_{ij})}{r_{ij}},
  \label{eq:sM}
\end{equation}
using relativistic DPWA (Dirac partial-wave analysis) electron
scattering factors\cite{Ross1992} for C, H, N, O on the grid
$s = 0.022$--$15.0\anginv$ (680 points). The training signal is
\begin{equation}
  \dsm_{\rm train}(s) = sM(s;\,\text{struct}) - sM(s;\,\text{equil}).
\end{equation}

\subsection*{C.\; Noise Model}

Experimental noise was characterized from the 200 bootstrap-resampled
pre-time-zero frames. The root-mean-square variation across bootstrap
samples and pre-pump time points in $s \in [1.5,\,9.0]\anginv$ was
$\sigma_{\rm exp}\approx 0.003$ (experimental units), corresponding
to $\approx 5.5$ in calibrated units after multiplication by $\alpha$.
Training samples are augmented with Gaussian white noise at standard
deviation drawn uniformly from $[0.002,\,0.02]$, plus a slow-drift
component (std~$\in$~$[0,\,0.01]$, smoothed over 31 grid points) to
simulate low-frequency residuals.

\subsection*{D.\; Network Architecture}

The regression network (NMMRegressorV2,
Fig.~\ref{fig:network}) has ${\sim}2.13\times10^6$ parameters:

\noindent\textit{Multi-scale stem.} Four parallel Conv1d branches
with kernel widths $k \in \{3,7,15,31\}$ capture oscillatory features
at different spatial frequencies, corresponding via the Fourier
relation to atomic pair correlations at different real-space distances.
Outputs are concatenated and projected to $C = 128$ channels.

\noindent\textit{Residual body.} Eight residual blocks\cite{He2016}
with Conv1d--BatchNorm--ReLU sequence and skip connections, arranged
in pairs. MaxPool1d downsamples after every two blocks (16$\times$
total). Each block includes a Squeeze-and-Excitation (SE)
attention\cite{Hu2018} module.

\noindent\textit{Regression head.} Global average pooling $\to$
FC(512) $\to$ FC(256) $\to$ linear(7).

\subsection*{E.\; Label Design}

The network predicts seven interatomic distances:
$d_{O\text{-}N}$, $d_{O\text{-}C5}$, $d_{N\text{-}C5}$ (three key
distances involving the two mobile atoms and oxygen), plus
$d_{N\text{-}C2}$, $d_{N\text{-}C4}$, $d_{C5\text{-}C2}$,
$d_{C5\text{-}C4}$ (four ancillary distances to the fixed anchor
atoms C2 and C4). Seven distances provide an over-determined system
for the six coordinates of N3 and C5, enabling unique 3D
reconstruction via trilateration (Section~IV) and serving as an
internal self-consistency check.

The three-distance label used in an earlier model phase was found to
be geometrically ambiguous: three distances from a common set of
anchors do not uniquely fix two points in 3D space. This led to
unreliable N--C5 predictions on experimental data despite low
validation MAE.

\subsection*{F.\; Training Protocol}

The AdamW optimizer\cite{Loshchilov2019}
($\eta_0 = 1.5\times10^{-4}$, weight decay $10^{-4}$) with 3-epoch
linear warmup, cosine annealing to zero at epoch~50,
dimension-weighted SmoothL1 (Huber)
loss with weights $\mathbf{w} = [1.0,\,1.0,\,2.5,\,1.5,\,1.5,\,3.0,\,3.0]$
(emphasizing C5--C2 and C5--C4 to improve C5 localization), and
gradient clipping (max norm~1.0) was used. Batch size: 128.
Hardware: Apple M3 GPU (MPS backend). An early-stopping criterion
(patience~=~15) on validation MAE was applied.

\subsection*{G.\; Dataset}

$8\times10^4$ training and $8\times10^3$ validation samples were
generated with \texttt{equil\_fraction}~=~0.3, meaning 30\% of
samples are drawn from a narrow Gaussian distribution around the
equilibrium geometry (Franck-Condon region), while the remaining 70\%
are uniformly sampled from the full cylinder$\times$sphere domain.
Each sample consists of a 681-point $\dsm_{\rm train}(s)$ vector
(input, normalized channel-wise) and the seven-distance vector
(output, in \AA).

%%=============================================================
\section*{IV.\; Results}
%%=============================================================

\subsection*{A.\; Synthetic Validation}

Table~\ref{tab:val} summarizes the validation-set MAE of the
best model (7-dim, 80K, \texttt{equil\_fraction}~=~0.3).
For reference, an earlier 3-dim model achieved key MAE =
0.081~\AA\ but suffered from geometric ambiguity and unreliable
N--C5 predictions on experimental data. The 7-dim model provides
both accurate distance prediction and unique 3D reconstructability.

\begin{table}[htb]
  \caption{Validation-set MAE of NMMRegressorV2 (7-dim, 80K,
    equil\_fraction=0.3) on all seven predicted distances.}
  \label{tab:val}
  \centering
  \begin{tabular}{lcc}
    \toprule
    Distance & MAE (\AA) & $R^2$ \\
    \midrule
    O$\cdots$N   & 0.035 & 0.993 \\
    O$\cdots$C5  & 0.065 & 0.986 \\
    N--C5        & 0.190 & 0.956 \\
    N--C2        & 0.146 & 0.957 \\
    N--C4        & 0.154 & 0.934 \\
    C5--C2       & 0.225 & 0.830 \\
    C5--C4       & 0.221 & 0.834 \\
    \midrule
    \textit{Key MAE (top 3)} & \textit{0.097} & \\
    \textit{Overall MAE}     & \textit{0.148} & \\
    \bottomrule
  \end{tabular}
\end{table}

\subsection*{B.\; Pre-time-zero Baseline}

The network applied to the 9 pre-pump frames ($t < 0$) with the
optimized calibration factor $\alpha = 0.02$ yields a mean deviation
of 5.7\% from the DFT equilibrium distances across all seven
dimensions. The two key distances O$\cdots$N and O$\cdots$C5 show
excellent agreement: 0.2\% and 1.2\% deviation, respectively. A
zero-input test (feeding an all-zero signal) produces distance
predictions within $\pm 3.4\%$ of equilibrium values, confirming
that the network has learned a physically reasonable prior for the
ground-state geometry. The calibration factor $\alpha = 0.02$ was
selected by grid search over $\alpha \in [0.02,\,0.04]$
(Fig.~S1), yielding the minimum mean pre-pump deviation.

\subsection*{C.\; Time-Resolved Distance Trajectories}

Figure~\ref{fig:trajectories} shows the bootstrap mean $\pm 1\sigma$
trajectories for the seven predicted distances over all 45 time delays.
Key observations:
\begin{itemize}
  \item \textbf{Pre-pump baseline} ($t < 0$, t00--t08): All
    distances are approximately consistent with ground-state values,
    $d_{O\text{-}N}^{\rm eq} = 2.851$\AAng,
    $d_{O\text{-}C5}^{\rm eq} = 4.222$\AAng,
    $d_{N\text{-}C5}^{\rm eq} = 1.455$\AAng, with an optimized
    calibration factor $\alpha = 0.02$ yielding mean pre-pump
    deviation of $\sim$5.7\% across all seven distances.
  \item \textbf{Rise dynamics} ($t \approx 0$--200~fs): The
    O$\cdots$N distance increases from 2.86~\AA\ (pre-pump) to a
    maximum of 3.31~\AA\ at $t \approx 134$~fs ($\Delta d = +0.45$~\AA,
    $+16\%$). The N--C2 and N--C4 distances also increase by $+0.37$
    and $+0.31$~\AA, respectively, consistent with nitrogen moving away
    from the ring plane toward a planar geometry.
  \item \textbf{O$\cdots$C5 contraction}: The O$\cdots$C5 distance
    decreases from 4.27~\AA\ to $\sim$3.86~\AA\ post-pump ($-$0.31~\AA),
    consistent with C5 being pulled closer to the ring as N3
    planarizes.
  \item \textbf{Excited-state evolution} ($t > 200$~fs):
    Distance trajectories show continued temporal variation,
    with the O$\cdots$N distance remaining elevated above the
    equilibrium value throughout the measured time window.
\end{itemize}

\subsection*{D.\; 3D Structure Reconstruction}

The seven predicted distances are sufficient to uniquely reconstruct
the 3D coordinates of N3 and C5 via trilateration from the three
fixed anchor atoms O7, C2, C4. The reconstruction algorithm
(i)~solves two independent trilateration problems for N3 and C5
(each yielding a mirror pair), then (ii)~selects the combination
minimizing $|d_{NC5}^{\rm recon} - d_{NC5}^{\rm pred}|$.

All 45 time points yield successful reconstructions (100\% success
rate), with self-consistency residuals
$|d_{\rm recon} - d_{\rm input}| < 0.01$~\AA\ for the
anchor-referenced distances. Pre-pump N3 displacement from
equilibrium is 0.34~\AA, and C5 displacement is 0.74~\AA\
(Fig.~\ref{fig:reconstruction}), reflecting residual domain bias
rather than physical motion.

%%=============================================================
\section*{V.\; Discussion}
%%=============================================================

\subsection*{A.\; Physical Interpretation}

Upon $\sim$200~nm photoexcitation, NMM enters the 3s Rydberg
state.\cite{Zhang2017,Beedle2017} The dominant nuclear response is the
nitrogen umbrella mode: N3 moves from the sp$^3$-pyramidal toward a
planar geometry, shortening the lone-pair extent along the ring axis.
This motion directly modulates the O$\cdots$N, N--C5, and O$\cdots$C5
distances accessible to our ML model. The extracted O$\cdots$N distance increases by $+0.45$~\AA\
within $\sim$130~fs, directly quantifying the nitrogen umbrella
inversion in physical distance units. This corresponds to an
N3 displacement of $\sim$0.3--0.5~\AA\ out of the ring plane,
consistent with the spectroscopically inferred 17° Franck-Condon
displacement\cite{Zhang2017} (which corresponds to $\sim$0.4~\AA\
for the $\sim$1.46~\AA\ N--C bond length). The O$\cdots$C5
contraction ($-0.31$~\AA) and N--C2/N--C4 elongation
($+0.3$--$0.4$~\AA) are consistent with the geometric consequences
of nitrogen planarization: as N3 moves into the ring plane, it
moves closer to the ring center and further from the C2/C4 ring
carbons. Future work with improved domain-bias mitigation may
reveal the expected $\sim$650~fs oscillatory dynamics.

\subsection*{B.\; Advantages and Limitations of ML Inversion}

Compared with traditional iterative structure refinement,\cite{Centurion2022}
the ML approach (i)~requires no prior knowledge of the reaction
coordinate, (ii)~runs in $<1$~ms per signal after training,
(iii)~propagates bootstrap uncertainty without rerunning an
optimization, and (iv)~is straightforwardly retrained for different
molecules by changing the synthetic training set. The primary
limitation is the domain gap between synthetic (IAM + DPWA) and
experimental signals (true electron density; structured detector
noise). The noise gap---the largest contributor to failure of earlier
low-noise models---is closed explicitly by matching training noise to
the experimental pre-pump standard deviation. Residual systematic
offsets from the IAM approximation warrant future investigation with
beyond-IAM scattering factors.\cite{Kirrander2021}

%%=============================================================
\section*{VI.\; Conclusion}
%%=============================================================

\noindent
We have developed and applied a noise-matched machine learning
inversion framework for MeV-UED data of N-methylmorpholine:
\begin{enumerate}
  \item NMMRegressorV2, trained on $8\times10^4$ synthetic DPWA-based $\dsm$
    signals with experimentally matched noise, achieves 0.097~\AA\
    key MAE and 0.148~\AA\ overall MAE on seven interatomic distances.
  \item Matching training noise to the experimental level closes the
    primary domain gap that caused failure of low-noise models, without
    additional model complexity.
  \item Bootstrap uncertainty propagation provides 1$\sigma$ distance
    uncertainties at each time point at negligible computational cost.
  \item The method reveals a $+0.45$~\AA\ ($16\%$) increase in
    O$\cdots$N distance within $\sim$130~fs of photoexcitation,
    directly quantifying the nitrogen umbrella inversion amplitude
    and providing the first direct structural measurement of this
    fundamental photochemical coordinate in NMM.
\end{enumerate}
The framework is directly applicable to other nitrogen and oxygen
heterocycles studied by gas-phase UED, offering a general path toward
real-time structural tracking of photochemistry in medium-sized
organic molecules.

%%=============================================================
\section*{Acknowledgments}

\TBD{Funding acknowledgements.}
The authors thank \TBD{collaborators} for helpful discussions.
UED experiments were performed at the MeV-UED facility of Shanghai
Jiao Tong University.

\section*{Data Availability}

The training code, trained model weights, and experimental data
supporting the conclusions of this study are available from the
corresponding author upon reasonable request.

%%=============================================================
\clearpage
\bibliographystyle{sciencemag}
\bibliography{paper_nmm}

%%=============================================================
%% FIGURES
%%=============================================================
\makeatletter
\renewcommand{\fnum@figure}{\textbf{Fig.\ \thefigure}}
\renewcommand{\fnum@table}{\textbf{Table\ \thetable}}
\makeatother

\clearpage
\section*{Figures}

\begin{figure}[h]
  \centering
  \fbox{\parbox{0.75\textwidth}{\centering\vspace{2.5cm}
    \textit{[Fig.~1 — NMM structure and sampling regions]}\vspace{2.5cm}}}
  \caption{
    \textbf{Ground-state structure of NMM and structural sampling regions.}
    (a)~Ball-and-stick representation (B3LYP/6-311++G**). Fixed atoms
    (C1, C2, C4, C6, O7) are shown in blue; mobile atoms (N3, C5) in
    orange. H atoms (gray) are rigidly coupled to parent heavy atoms.
    (b)~Sampling regions: N3 cylinder (radius 2.0~\AA, off-plane
    2.0~\AA) and C5 sphere (radius 3.0~\AA\ centered on ground-state
    C5). Gray shading indicates the acceptance region after the
    constraint $r_{N\text{-}O} < r_{C5\text{-}O}$.
    (c)~Seven interatomic distances predicted by the network:
    three key distances ($d_{O\text{-}N}$, $d_{O\text{-}C5}$,
    $d_{N\text{-}C5}$) and four ancillary distances to C2/C4 anchors.
  }
  \label{fig:structure}
\end{figure}

\begin{figure}[h]
  \centering
  \fbox{\parbox{0.75\textwidth}{\centering\vspace{2.5cm}
    \textit{[Fig.~2 — Network architecture]}\vspace{2.5cm}}}
  \caption{
    \textbf{NMMRegressorV2 architecture.}
    Input: 681-point $\dsm(s)$ signal (normalized).
    \textit{Multi-scale stem}: four parallel Conv1d branches with
    kernel widths 3, 7, 15, 31, capturing reciprocal-space oscillations
    corresponding to atomic pair correlations at different length scales.
    \textit{Residual body}: 8 ResBlocks with SE channel attention and
    MaxPool1d downsampling (2 pooling steps, $4\times$ total).
    \textit{Regression head}: global average pooling $\to$
    FC(512) $\to$ FC(256) $\to$ linear(7).
    Total parameters: ${\sim}2.13\times10^6$.
  }
  \label{fig:network}
\end{figure}

\begin{figure}[h]
  \centering
  \includegraphics[width=0.95\textwidth]{Figures/fig3_training_curves.png}
  \caption{
    \textbf{Training progress across model generations.}
    (a)~Overall validation MAE as a function of training epoch for
    the 3-dim 80K model (blue), 7-dim 80K model (red), and 7-dim
    80K model with equilibrium fraction (green, current best).
    Dots indicate the best epoch for each model.
    (b)~Per-distance validation MAE vs.\ epoch for the 7-dim
    equilibrium model, showing convergence behavior for all seven
    predicted distances.
  }
  \label{fig:noise}
\end{figure}

\begin{figure}[h]
  \centering
  \includegraphics[width=0.95\textwidth]{Figures/fig4_parity.png}
  \caption{
    \textbf{Network validation on synthetic data.}
    (a)--(g)~Parity plots (predicted vs.\ true) for all seven
    distances on the 8,000-sample validation set.
    The diagonal dashed line indicates perfect prediction; MAE and
    $R^2$ values are given in each panel.
    (h)~Learning curve: validation key-MAE vs.\ training epoch.
  }
  \label{fig:validation}
\end{figure}

\begin{figure}[h]
  \centering
  \includegraphics[width=0.85\textwidth]{Figures/fig5_trajectories.png}
  \caption{
    \textbf{Time-resolved interatomic distances from ML inversion.}
    Bootstrap mean (solid lines) $\pm 1\sigma$ uncertainty (shaded
    bands, $N = 200$ resamples) for all seven predicted distances as
    functions of pump-probe delay. Dashed horizontal lines: DFT
    equilibrium values. Vertical dashed line: time zero. Gray region
    ($t < 0$): pre-pump reference window.
  }
  \label{fig:trajectories}
\end{figure}

\begin{figure}[h]
  \centering
  \includegraphics[width=0.95\textwidth]{Figures/fig6_reconstruction.png}
  \caption{
    \textbf{3D molecular structure reconstructed from predicted distances.}
    Six time-point snapshots showing the positions of N3 (blue) and C5
    (green) reconstructed via trilateration from the seven predicted
    distances. Fixed ring atoms (gray) and O7 (red) are at ground-state
    positions. Cross markers indicate equilibrium positions. Bond lengths
    (N--C5 and O$\cdots$N) are annotated in each frame.
  }
  \label{fig:reconstruction}
\end{figure}

\end{document}
